\section{Introduction}
A finite future experiment tests only part of a posterior distribution.
The amount of information needed to predict it is determined neither by
the dimension of that test space alone nor by a separate quantizer at
each time. The past must actually acquire the relevant directions; those
directions can become indistinguishable as the experiment varies; and
one retained state must support all subsequent causal updates. If the
observation order is itself selected from data, the probabilities of
visiting different information states must also be compatible across
checkpoints.

We study these constraints through attainable prediction geometry and
common occupation measures. The first principal result is a
collision-uniform classification for positive sparse monomial
experiments, at every integer label budget. The second transports the
acquired geometry to adaptive graph experiments and distinguishes an
order chosen at one resolution from an acquisition law reused over a
range of resolutions. The third identifies which conclusions follow
from local path probabilities and which require additional
nonanticipativity information. On regenerative acquisition networks,
that additional information yields an exact statistical minimax formula,
a causal realization, and a controlled gap between small-ball flow
bounds and the actual optimum. The proofs separate these hypotheses
rather than identifying an abstract relaxation with an arbitrary
statistical prediction problem.

\subsection{Attainable information and collision-uniform resolution}
Let
\[
 A=\{0=a_0<a_1<\cdots<a_{r-1}\},\qquad
 W_A=\operatorname{span}\{t^{a_i}:0\le i<r\},\qquad r\ge2.
\]
We put $I=[l,u]=[0,1]$, $D=a_{r-1}$, and
\[
 mA=\{a_{i_1}+\cdots+a_{i_m}:0\le i_j<r\},\qquad
 h_A(m)=|mA|,
\]
where $0A=\{0\}$ and $h_A(0)=1$. The set $mA$ contains distinct
values; the formal labels used below retain multiplicities. The latent
prior $\mu$ is a fixed full-support probability on $I$. It need not
have a density.

The known exponent vector ranges over a compact subset $K$ of the
strictly ordered one-step chamber. A fixed matrix $C=(c_{ji})$ of
column rank $r$ defines detector cells
\[
 k_j^a(t)=\sum_i c_{ji}t^{a_i},\qquad
 \sum_j k_j^a(t)=1,\qquad k_j^a(t)\ge k_*>0
                  \quad(a\in K,\ t\in I).
\]
An exogenous command $g\in[\eta,1-\eta]^J$, with $0<\eta<1/2$,
retains cell $j$ with probability $g_j$. A trial reports that cell or
failure. Every trial, including failure, is counted. The report
likelihoods are bounded below by one fixed $\kappa>0$.
Section~\ref{sec:experiment}
defines the complete experiment and its future-only continuations.

Fix a total horizon $N\ge2$ and a checkpoint $n$, with $m=N-n$.
Choose $H\ge\max\{1,N\max_{a\in K}a_{r-1}\}$ and set
\[
 \Lambda_m^+=\{\alpha\in\N_0^r:|\alpha|=m\}\setminus\{me_0\},
 \quad q_m=\binom{m+r-1}{r-1}-1,
 \quad x_\alpha(a)=H^{-1}\alpha\cdot a.
\]
Here $\N_0$ denotes the nonnegative integers. Define the maximal
future determinant volumes
\[
 \mathcal V_{m,\ell}(a)=
 \max_{\substack{F\subset\Lambda_m^+\\|F|=\ell}}
       \prod_{\{\alpha,\beta\}\subset F}
                      |x_\alpha(a)-x_\beta(a)|,
                      \qquad \mathcal V_{m,1}=1.
\]
Zero products are allowed, and formal labels remain present at exact
collisions. The acquired truncation and resolution profiles are
\[
 \begin{split}
 p_{n,m}&=\min\{n(r-1),q_m\},\\
 \Psi_{n,m}(M,a)&=
   \max_{1\le\ell\le p_{n,m}}
       \left(\frac{\mathcal V_{m,\ell}(a)}{M}\right)^{2/\ell},\\
 \Xi_N(M,a)&=\max_{1\le n<N}\Psi_{n,N-n}(M,a).
 \end{split}
\]
Endpoint profiles are zero. The normalization $H$ affects only
comparison constants.

A checkpoint encoder reads the acquired prefix and retains one of $M$
labels. A causal filter updates only its previous label from the current
command and report. After compression, an independent query is selected
from a fixed finite spanning menu of actual failure-product
continuations. Only the selected continuation is executed. Excess
squared loss is measured against the full-history conditional mean.
Definitions~\ref{def:finite-state} and
\ref{def:monomial-risks-v24} specify the retained resource and the
order of the risk quantifiers.

\begin{theorem}[Intrinsic attainable resolution and causal memory]\label{thm:resolution-main}
Under these hypotheses there are constants $0<c\le C<\infty$,
depending only on the fixed experiment, $K$, $N$ and $\mu$, with the
following properties for every known calibration $a\in K$ and every
integer $M\ge1$.

At each checkpoint $n+m=N$, the optimal $M$-message squared-prediction
regret is between $c\Psi_{n,m}(M,a)$ and $C\Psi_{n,m}(M,a)$,
both unconditionally under independent uniform acquisition commands and
in the minimax sense over admitted histories. The query is sampled
independently after the index is formed, from the ordered products of
one fixed attainable failure-factor basis; all its $m$ trials are counted.

There exists one deterministic stage-compatible transducer, allowed to
depend on $a$ and $M$, storing at most $M$ persistent labels after
every report, whose regret at every checkpoint and on every history is
at most $C\Xi_N(M,a)$. Conversely every such transducer has maximum
unconditional checkpoint regret at least $c\Xi_N(M,a)$ under the
same common exploration law. Independent coding randomization does not
remove the lower bound. The clock and fixed real-valued program are
read-only as in Definition~\ref{def:finite-state}.

The constants are uniform across all additive collision strata of $K$.
No calibration-blind codebook or uniformity over all full-support priors
is part of the assertion.
\end{theorem}


At a fixed calibration the last nonzero profile index is
\[
             d_A(n,m)=\min\{n(r-1),|mA|-1\}.
\]
Theorem~\ref{thm:main} identifies the exact continuous-state dimension.
The determinant volumes carry the information needed to compare
experiments uniformly through changes of that dimension. Theorem~
\ref{thm:resolution-main} is therefore an all-budget result, not only
an eventual error exponent at a separately fixed calibration.

\subsection{Intersecting collision loci}
\label{sec:early-collision-example}
Consider $A_{u,v}=\{0,1,2+u,3+v\}$ for $|u|,|v|\le1/16$, with
\[
 k_0(t)=\frac14+\frac{t+t^{2+u}+t^{3+v}}{16},\qquad
 k_i(t)=\frac14-\frac{t^{a_i}}{16}\quad(1\le i\le3).
\]
These cells have a fixed rank-four coefficient matrix and lie between
$3/16$ and $7/16$. For $N=5$, write
\[
 \rho=\max\{|u|,|v|\},\qquad
 \tau=\min\{|u|,|v-u|,|v-2u|\}.
\]
Corollary~\ref{cor:two-parameter} gives both causal risk criteria the
uniform order
\begin{equation}\label{eq:early-two-parameter-v24}
 \max\left\{M^{-1/3},\ \rho^{1/2}M^{-1/4},\
                    (\rho^2\tau)^{2/9}M^{-2/9}\right\}.
\end{equation}
The peak acquired dimension is nine away from the three collision
lines, eight on a line away from their intersection, and six at the
intersection. At $(n,m)=(3,2)$ the nine positive formal sums form six
separated groups, with internal gaps $|u|$, $|v-u|$, and $|v-2u|$.
The largest two are comparable, making the apparent seven-dimensional
branch redundant. This is a property of the geometry, not a missing
term. Along $u=\theta$, $v=\theta+\theta^k$, $k\ge2$, the crossover
budget orders are $\theta^{-6}$ and $\theta^{-(8k-2)}$, and the
corresponding regret orders are $\theta^2$ and $\theta^{2k}$.
Their multiplicative constants are not claimed to be exact switching
thresholds. The general contact-order result is
Theorem~\ref{thm:collision-tree}.

\subsection{One acquisition law and one causal state}
\label{sec:contribution-v23}\label{sec:proof-route-v24}
The acquired truncation is substantive. For $A=\{0,1,3\}$ and
$(n,m)=(1,2)$ there are five nonconstant future exponents, but only
two acquired directions; the eventual squared-error order is $M^{-1}$,
not $M^{-2/5}$. The rank is obtained from executable commands: the
product tangent at a binomial history has dimension $n(r-1)+1$,
strict mixed-moment positivity gives its pairing rank, and evidence
normalization removes one direction.

Uniformity through collisions requires complete confluent flags rather
than isolated derivatives. A finite Leja ordering supplies decreasing
Newton scales whose prefix products compare with maximal determinant
volumes. The same binomial tangent pairs with every complete Hermite
prefix. An inverse chart includes complementary command coordinates;
integrating them, while retaining the actual report evidence, gives
positive acquired mass under one unconditional exploration law
(Lemma~\ref{lem:newton-attainment}). No density assumption on the
latent prior is substituted for this acquired-law statement.

For the upper bound, each fixed-word attainable image is rational in
its command variables with positive denominator. A bounded-format
cover controls the entire image, not only the lower-bound chart.
Theorem~\ref{thm:intrinsic-checkpoint} combines this cover with the
acquired-prefix lower measures. The raw remaining-moment recursion
has positive Bayes denominators and no reciprocal collision gap.
Reachable representatives can therefore be updated and re-quantized,
retaining every preceding approximation error. This gives one filter
for all checkpoints (Theorem~\ref{thm:intrinsic-streaming}); the
converse places its infimum outside the checkpoint maximum under the
same exploration law.

\subsection{Adaptive graphs and reusable acquisition}
Put independent scalar experiments on the edges of a fixed graph.
Visiting the first endpoint acquires an edge block; visiting its second
endpoint completes it. Memory is charged at vertex-batch boundaries,
and the future query is the full tensor product of the remaining probe
menus. All raw trials are counted. Each edge contributes a capped
profile
\[
 b_e(\varepsilon,a)=\max_{0\le\ell\le p_e}
       \left\{\frac{\ell}{2}\log_2\frac1\varepsilon
                           +\log_2V_{e,\ell}(a)\right\},
\]
with zero branch zero and zero positive volumes omitted. Theorem~
\ref{thm:v25-adaptive} and Corollary~\ref{cor:v25-graph-bits} identify
the persistent-bit law, to a calibration-uniform bounded additive term,
as the minimum over vertex orders of the maximum cut sum of these
weights. The adaptive converse allows measurable history-dependent
selection and retains unconditional product-tape measures before
restriction to a selected trace. The scalar acquired flags, rather than
an assumed ambient product dimension, determine the edge caps.

The order minimizing this profile may depend on resolution. A
resolution-blind rule instead has fixed acquisition transitions and a
fixed independent seed law across predictive budgets; it receives no
feedback from a budget-dependent encoder. It may read the full revealed
history in the stated converse relaxation. Theorem~\ref{thm:v27-curve}
selects one deterministic comparison order for its entire prediction
curve. Theorem~\ref{thm:v27-menu} extends this to a finite menu whose
index is chosen before data. The predictive codebooks, unlike the
acquisition law, may be redesigned at every budget.

At every fixed graph and calibration, one order is eventually optimal
and has finite excess over all accuracies in $(0,1]$
(Theorem~\ref{thm:v30-menu-finite}). Proposition~
\ref{prop:v30-critical-resolutions} expresses the optimized all-resolution
menu excess as a finite critical-resolution minimax problem. This does
not make the excess uniformly bounded as calibration approaches a
collision. In the actual three-edge, $24$-trial star, the optimal
single-rule excess over the indicated resolution interval is
$\log_2(1/\theta)+O(1)$, while two configurations remove its leading
term (Corollary~\ref{cor:v27-two-configurations}). The fixed-calibration
and collision-uniform quantifiers are different.

The size-uniform graph theorem has its own loss convention. Theorem~
\ref{thm:v28-localized} uses selected-edge witnesses at a compulsory
boundary and retains local density and recovery constants. Under the
strong retained balanced-cut and independence hypotheses, Theorem~
\ref{thm:v28-size-uniform} gives an extensive bit law for linear-size
graph families. Its menu gives half the query mass to original tensor
probes and half to local edge readouts. Equality of their exact test
spaces is not uniform equivalence of these weighted loss metrics as
the graph grows. No tensor-only lower bound is inferred from the
augmented-task theorem. Its nonempty graph families and joint
collision-growth consequence are established in
Proposition~\ref{prop:v28-graph-existence} and
Corollary~\ref{cor:v28-joint-limit}.

\subsection{From local path bounds to statistical realization}
For a layered acquisition graph, suppose $c_w(t)$ bounds the normalized
probability of visiting vertex $w$ with loss at most $t$ on one event
of mass $\beta$. Put $\phi_w(x)=\int_0^1[x-c_w(t)]_+\,dt$.
The occupation converse is
\[
 \max_j\E L_j\ \ge\ \beta
    \min_{x\in\mathcal F(c)}\max_j
                   \sum_{w\in\mathscr V_j}\phi_w(x_w).
\]
The same flow is used at every threshold and level, with its terminal
node caps retained. Theorems~\ref{thm:v29-occupation} and
\ref{thm:v29-relaxation-exact} show that this is exactly the strongest
consequence of the stated local path data. Conditioning on independent
seeds before the first acquisition yields a stronger perspective bound
and an exact dual (Theorems~\ref{thm:v29-pre} and \ref{thm:v29-dual}).
A feasible dual gives a lower certificate; a nonoptimal primal does not.
These exactness assertions concern their specified relaxations, not
arbitrary posterior or decoder realizability.

Under conditional regeneration, one can go further. A vertex is chosen
from the revealed history before its fresh acquisition block is
observed; given that choice and the past, the new block has its prescribed
law. The decoder receives the vertex as its finite query descriptor
and one of $M$ labels, but no discarded history. If $d_w(M)$ is the
one-block optimal distortion at vertex $w$, then
Theorem~\ref{thm:v30-predictable-exact} proves the exact statistical
identity
\begin{equation}\label{eq:v30-intro-network}
 \mathcal R_{\mathscr D,M}^{\rm reg}
   =\min_{x\in\mathcal F}\max_j\sum_{w\in\mathscr V_j}x_wd_w(M)
   =\max_{\lambda\in\Delta_J}\min_{P:o\rightsquigarrow d}
                       \sum_{w\in P}\lambda_{j(w)}d_w(M).
\end{equation}
A distribution on at most $J+1$ deterministic paths and a single causal
reset quantizer attain the optimum. The lower bound even permits a
full-history scheduler. Thus the flow is realized by statistical
acquisition, finite decoding and a common causal controller, rather
than by independently assigned losses at the network vertices.

For positive two-trial blocks with the full acquired mixed law $\nu$
and known positive readout gains $s_w$, the one-block distortion is
$s_w^2D_M(\nu)/128$. The exact network value factors into this scalar
distortion and one common design value. Predictable atom-aware
small-ball bounds give a certificate for which the ratio of the true
optimum to the certificate is at most $8306688/64800<129$ at every budget and tends to a number
in $(4.05,4.06)$, uniformly over all finite network sizes, horizons and
positive gains (Corollary~\ref{cor:v30-network-gap}). The full acquired
law is used at each stage, so there is no exponentially small event
requiring every stage to report failure.

The common flow cannot be replaced by separate checkpoint optimizers.
A realized $J$-route, $J$-checkpoint family has an exact ratio
$J+1-1/J$ between its true optimum and that invalid separate-checkpoint
comparison at every finite budget
(Corollary~\ref{cor:v30-checkpoint-gap}). These examples use an explicit
regenerative and attenuated readout protocol. Conditional regeneration
is not asserted for the preceding nonregenerative graph posterior, and
the public vertex descriptor is part of the protocol, with its optional
storage cost stated separately. This structural result complements the
collision classification; it does not replace its experiment by a
simpler one.

\subsection{Classical inputs and organization}
Newton and Hermite interpolation, Leja orderings, and bounded-format
covering are classical tools
\cite{deBoor,Leja,BosEtAl,YomdinComte,ZhangKileel,ComteHalupczok}.
They do not by themselves provide positive acquired mass in an
experiment or a converse for a common finite-label filter. Nonlinear
filter quantization and error propagation are developed in
\cite{PagesPham2005,PagesSagna2018}. An approximation with an
unrestricted vector of filtering weights has a different resource
from the complete retained prediction label counted here; the companion
comparison makes that distinction at the level of mathematical objects.
Scalar high-resolution quantization is classical
\cite{GrafLuschgy2000}, as are finite convex separation and minimax
\cite{Rockafellar1964}. No originality is attributed to those tools.
The statistical claims above include the acquired law, the permitted
information, and the compatibility of a common causal implementation.

Sections~\ref{sec:experiment}--\ref{sec:collision-phases} establish
the scalar classification and its consequences. The graph sections
then develop adaptive selection, reusable menus, local witnesses and
growing experiments. Sections~\ref{sec:v29-occupation} and
\ref{sec:v29-initial} separate local-path exactness from initial
information. Sections~\ref{sec:v30-multilevel} and
\ref{sec:v30-predictable} give statistical realization and predictable
network duality. The fully evaluated acquired law and its scalar
sharp constant are proved in Section~\ref{sec:v27-evaluated}, with
its complete earlier quantitative comparisons retained in
Section~\ref{sec:v29-evaluated}.

The accompanying volume contains the complete transfer, bounded
observation-algebra, exact-kernel, uncertainty, operational and
effective-construction developments. Its sections have prefix C;
references back to this article have prefix M. The bounded-algebra
classification is uniform over priors and support losses but has
immediate acquisition saturation. It is a separate result, not an
additional hypothesis or an all-prior interpretation of the monomial
theorem. In that theorem the prior and horizon are fixed, and its
constants are uniform in calibration and label budget. All historical
proof modules remain available; the organization here follows
mathematical dependence rather than the sequence of revisions.
